\section{Colored Tornheim functions and entire continuation}\label{tsc:tornheim}
For an initially convergent double sum, write
\begin{equation}\label{tsc:Tdef}
 T(r,s,t;z,w)=\sum_{m,n\ge1}
           \frac{z^m w^n}{m^r n^s(m+n)^t}.
\end{equation}
For example, $\Re r,\Re s,\Re t>1$ is more than sufficient for unit
colors. The functions below are not defined by assigning a value to a
divergent instance of this series. They are defined by its unique analytic
continuation, constructed explicitly in the next theorem.

\begin{theorem}[Nontrivial-color entireness]\label{tsc:T-entire}
Fix $|z|=|w|=1$ with $z\ne1$ and $w\ne1$. Then
$T(r,s,t;z,w)$ extends to a jointly entire function of $(r,s,t)\in\CC^3$.
For $\Re t>0$ and arbitrary $r,s\in\CC$ it is given by
\begin{equation}\label{tsc:T-mellin}
 T(r,s,t;z,w)=\frac1{\Gamma(t)}
   \int_0^\infty x^{t-1}\Li_r(ze^{-x})\Li_s(we^{-x})\dd x.
\end{equation}
The two colors may coincide. At every nonnegative integer $k$,
\begin{equation}\label{tsc:T-negative}
 T(r,s,-k;z,w)=\sum_{j=0}^k\binom{k}{j}
                \Li_{r-j}(z)\Li_{s-k+j}(w).
\end{equation}
\end{theorem}
\begin{proof}
Put $F(x)=\Li_r(ze^{-x})\Li_s(we^{-x})$. Because neither color is one,
the arguments have neighborhoods avoiding the polylogarithm singularity
at one. The standard analytic continuation of $\Li_r$ is entire in its
order there \cite{tsc:dlmf-poly}. Thus $F$ is analytic in $x$ near zero,
jointly with $r,s$. On compact sets of orders, the power series for large
$x$ gives an exponentially decreasing bound for $F$ and its spectral
derivatives. These facts prove convergence and holomorphy of
\eqref{tsc:T-mellin} for $\Re t>0$. Expanding the two polylogarithms where
absolute convergence permits and integrating $e^{-(m+n)x}$ identifies it
with \eqref{tsc:Tdef} on a nonempty open set.

The Taylor coefficients at the lower endpoint are
\begin{equation}\label{tsc:Fcoeff}
 f_k=[x^k]F(x)=\frac{(-1)^k}{k!}
        \sum_{j=0}^k\binom{k}{j}
                   \Li_{r-j}(z)\Li_{s-k+j}(w).
\end{equation}
For any positive integer $N$, the expression
\begin{align}\label{tsc:T-subtracted}
 \frac1{\Gamma(t)}\bigg\{&
  \int_0^1 x^{t-1}\left(F(x)-\sum_{k=0}^{N-1}f_kx^k\right)\dd x
  +\int_1^\infty x^{t-1}F(x)\dd x
  +\sum_{k=0}^{N-1}\frac{f_k}{t+k}\bigg\}
\end{align}
is holomorphic for $\Re t>-N$ after removing the apparent singularities
at $0,-1,\ldots,-N+1$. Local uniformity in $r,s,t$ follows from the Taylor
remainder and exponential decay. The only poles of the braces are the
listed simple ones, and the simple zeros of $1/\Gamma(t)$ cancel them.
The expressions for different $N$ agree on overlaps, which proves joint
entireness. At $t=-k$, the zero of $1/\Gamma(t)$ has leading coefficient
$(-1)^k k!$; multiplication by the residue $f_k$ gives
\eqref{tsc:T-negative}.
\end{proof}

\begin{remark}[Why the color assumption matters]
At $z=1$ or $w=1$, the Taylor argument at $x=0$ changes: a polylogarithm
can have a fractional-power or logarithmic local term. The theorem does
not assert global entireness in that case. Later primitive formulas can
allow color one near positive spectral integers because the original
Tornheim series then converges absolutely. These are different reasons
for analyticity, and they must not be conflated.
\end{remark}

\subsection{An explicit jet at the zero third order}
For fixed nontrivial colors, define ordinary convergent logarithmic moments
\begin{align}\label{tsc:Mmoments}
 M_j(r,s;z,w)={}&\int_0^1\frac{\log^j x}{x}
        \bigl(\Li_r(ze^{-x})\Li_s(we^{-x})-\Li_r(z)\Li_s(w)\bigr)\dd x\notag\\
 &+\int_1^\infty\frac{\log^j x}{x}
                    \Li_r(ze^{-x})\Li_s(we^{-x})\dd x.
\end{align}
Taking $N=1$ in \eqref{tsc:T-subtracted} yields, locally for $|t|<1$,
\begin{equation}\label{tsc:T-zero-jet}
 T(r,s,t;z,w)=\frac1{\Gamma(1+t)}
  \left\{\Li_r(z)\Li_s(w)+
               \sum_{j\ge0}\frac{M_j(r,s;z,w)}{j!}t^{j+1}\right\}.
\end{equation}
Every derivative in $r,s$ may be passed through the moments. In particular,
\[
 T(r,s,0;z,w)=\Li_r(z)\Li_s(w),\qquad
 T_t(r,s,0;z,w)=\gamma\Li_r(z)\Li_s(w)+M_0(r,s;z,w).
\]
These formulas specify the third-order jets by convergent integrals rather
than by a limiting procedure at a pole.

\subsection{Classical double-polylogarithm and harmonic coordinates}
Use the convention
\[
 \Li_{p,q}(z,w)=\sum_{N>n\ge1}\frac{z^Nw^n}{N^p n^q}
\]
where convergent, followed by analytic continuation when justified.
Changing variables $N=m+n$ gives
\begin{equation}\label{tsc:T-zero-first}
 T(0,s,t;z,w)=\Li_{t,s}(z,w/z).
\end{equation}
For $z\ne w$, summing the finite geometric progression at fixed $m+n$
also gives
\begin{equation}\label{tsc:T00}
 T(0,0,t;z,w)=\frac{w\Li_t(z)-z\Li_t(w)}{z-w}.
\end{equation}
For positive integers $r,s$, the standard partial-fraction reduction
underlying Zhao's result \cite{tsc:zhao} is
\begin{align}\label{tsc:integer-reduction}
 T(r,s,t;z,w)={}&\sum_{j=1}^r\binom{r+s-j-1}{s-1}
                     \Li_{t+r+s-j,j}(w,z/w)\notag\\
 &+\sum_{j=1}^s\binom{r+s-j-1}{r-1}
                     \Li_{t+r+s-j,j}(z,w/z).
\end{align}
For completeness, the rational identity being summed is
\[
 \frac1{m^r n^s}=
  \sum_{j=1}^r\frac{\binom{r+s-j-1}{s-1}}{m^j(m+n)^{r+s-j}}
 +\sum_{j=1}^s\frac{\binom{r+s-j-1}{r-1}}{n^j(m+n)^{r+s-j}}.
\]
It follows by repeated use of
$1/(mn)=1/(m(m+n))+1/(n(m+n))$, or by comparing principal parts in
$m$ for fixed $m+n$. Multiplying by $(m+n)^{-t}z^mw^n$ and summing proves
\eqref{tsc:integer-reduction} first in an absolutely convergent domain.

At unit arguments and in a convergence domain,
\[
 \Li_{p,q}(1,1)=\sum_{N\ge1}\frac{H_{N-1}^{(q)}}{N^p},
 \qquad H_N^{(q)}=\sum_{j=1}^N j^{-q}.
\]
Thus the colored coordinate class includes generalized harmonic Euler sums.
The integer reduction is a known identity, not a new theorem here.
Moreover, differentiating it in the freely complex variable $t$ is
legitimate; differentiating it in the discrete integers $r$ or $s$ is not
licensed by the formula. The whole point of Theorem~\ref{tsc:T-entire}
is to supply a genuine analytic family for those missing jets.
